\documentclass[11pt]{article}
\usepackage[a4paper,margin=1in]{geometry}
\usepackage{amsmath,amssymb,mathtools,bm}
\usepackage{enumitem,booktabs,array}
\usepackage{tcolorbox}
\usepackage{hyperref}
\usepackage[protrusion=true,expansion=false]{microtype}
\hypersetup{colorlinks=true,linkcolor=blue,urlcolor=blue,citecolor=blue}
\setlist[itemize]{topsep=2pt,itemsep=2pt}
\setlist[enumerate]{topsep=2pt,itemsep=2pt}
\newtcolorbox{keybox}{colback=blue!3!white,colframe=blue!40!black,boxrule=0.4pt,arc=1mm}
\newtcolorbox{alertbox}{colback=red!3!white,colframe=red!40!black,boxrule=0.4pt,arc=1mm}
\newtcolorbox{answerbox}{colback=green!3!white,colframe=green!40!black,boxrule=0.4pt,arc=1mm}
\newtcolorbox{drillbox}{colback=yellow!5!white,colframe=orange!60!black,boxrule=0.4pt,arc=1mm}
\newcommand{\EV}{\mathbb{E}}
\newcommand{\Var}{\mathrm{Var}}
\newcommand{\Cov}{\mathrm{Cov}}
\newcommand{\blank}[1]{\underline{\hspace{#1}}}

\title{W12-02: Cline, Walkling \& Yore (2018, JFE) \\
\large ``The Consequences of Managerial Indiscretions: Sex, Lies, and Firm Value'' --- Active Recall Workbook}
\author{Banking \& Risk Management --- Exam Prep}
\date{\today}
\begin{document}\maketitle

\begin{keybox}
\textbf{Paper in one sentence.} Public disclosures of personal indiscretions (affairs, substance abuse, dishonesty) by US executives lead to a $\sim$1.6\% announcement-day abnormal stock return drop, increased subsequent governance failures (fraud, class-action suits), and elevated CEO turnover --- personal misconduct is a distinct \emph{social risk} that the market prices.

\textbf{Placement.} Defines a ``social'' risk channel beyond CSR and environmental; complements Lins-Servaes-Tamayo (W12-01) and extends the ESG-pricing frame to individual executive behaviour.
\end{keybox}

\section*{Round 1: Complete Walk-Through}

\subsection*{1.1 Sample and event}
Hand-collected sample of 325 public disclosures of US executive personal indiscretions 1978--2012, including marital infidelity, sexual harassment, substance abuse, dishonesty. Event study around news date.

\subsection*{1.2 Event-study results}
\begin{itemize}
\item Mean 3-day cumulative abnormal return around disclosure: $-1.6\%$ (statistically significant).
\item Effect larger for CEO than for lower-ranked executive indiscretions.
\item Effect larger for younger firms, higher-profile firms.
\end{itemize}

\subsection*{1.3 Subsequent outcomes}
\begin{itemize}
\item \textbf{Fraud.} Probability of subsequent financial fraud rises by $\sim$14\% at indiscretion firms vs.\ matched controls.
\item \textbf{Class-action suits.} More shareholder class actions post-event.
\item \textbf{Turnover.} CEO turnover more likely; successors perform better.
\item \textbf{Governance reforms.} More board changes, independent-director additions.
\end{itemize}

\subsection*{1.4 Mechanism}
\begin{itemize}
\item Indiscretions are proxies for lax ethical standards $\Rightarrow$ governance failures follow.
\item Reputation damage harms customer/employee trust.
\item Board acts to restore reputation, often via executive change.
\end{itemize}

\subsection*{1.5 Caveats}
\begin{alertbox}
\begin{itemize}
\item Selection: only \emph{public} indiscretions, possibly correlated with firm visibility.
\item Reverse causality: struggling firms may draw more scrutiny to executive behaviour.
\item Sample is a convenience sample; external validity limited.
\end{itemize}
\end{alertbox}

\section*{Round 2: Level 1 Fill-in}
\begin{enumerate}
\item Sample: \blank{1cm} US executive indiscretions, \blank{1cm}--\blank{1cm}.
\item Event-study CAR: $\sim$\blank{1cm}\%.
\item Fraud probability change: $\sim$\blank{1cm}\% increase.
\item CEO turnover: \blank{1cm} (more/less likely).
\item Mechanism: indiscretions signal lax \blank{1cm} standards.
\end{enumerate}

\section*{Round 3: Level 2 Prompts}
\begin{enumerate}
\item Write the event-study abnormal-return specification and discuss biases.
\item Why is indiscretion predictive of later fraud?
\item Discuss the market's reaction vs.\ the economic magnitude of the underlying misconduct.
\item How does CWY extend the ESG-pricing literature?
\end{enumerate}

\section*{Round 4: Minimal Scaffolding}
From a blank page: (i) sample, (ii) event CAR, (iii) subsequent outcomes, (iv) mechanism, (v) caveats.

\section*{Round 5: Blank Page}
Essay: managerial indiscretions as a distinct social-risk channel.

\vspace{6pt}
\begin{drillbox}
\textbf{Speed Drill.} (1) Sample size/years. (2) Event CAR. (3) Fraud likelihood change. (4) Turnover direction. (5) Two subsequent outcomes.
\end{drillbox}

\section*{Quick Reference \& Answer Key}
\begin{answerbox}
\begin{itemize}
\item \textbf{Sample.} 325 US executive indiscretions, 1978--2012.
\item \textbf{CAR.} $-1.6\%$ on announcement.
\item \textbf{Follow-on.} More fraud, more class actions, higher turnover.
\item \textbf{Mechanism.} Indiscretion signals weak ethics $\Rightarrow$ governance failures.
\end{itemize}
\end{answerbox}

\begin{keybox}
\textbf{30-second exam pitch.} Cline, Walkling \& Yore (2018) identify a distinct \emph{social risk} channel tied to individual executive behaviour. Using 325 hand-collected public disclosures of US executive personal indiscretions 1978--2012, they show a mean 3-day CAR of $-1.6\%$ on announcement. Over subsequent years, indiscretion firms exhibit $\sim$14\% higher fraud probability, more shareholder class actions, and elevated CEO turnover, with board composition changes following. The interpretation is that personal indiscretions proxy for lax ethical standards that show up later in governance failures. The paper extends ESG pricing from firm-level environmental and social metrics to individual executive behaviour and shows the market prices such reputation risk in both announcement returns and subsequent real outcomes.
\end{keybox}

\end{document}
